\documentclass[10pt]{article}
\usepackage[a4paper,margin=0.72in]{geometry}
\usepackage{amsmath,amssymb,amsthm,mathtools}
\usepackage{graphicx}
\usepackage[hidelinks]{hyperref}
\usepackage{xcolor}
\usepackage{microtype}
\setlength{\parindent}{0pt}
\setlength{\parskip}{5pt}
\newtheorem{theorem}{Theorem}
\newtheorem{lemma}{Lemma}
\newcommand{\Lp}{L^p}
\newcommand{\J}{\mathcal J}
\newcommand{\F}{\mathcal F}
\newcommand{\M}{\mathcal M}
\newcommand{\N}{\mathcal N}

\title{A piecewise-homogeneous tracial case of the positive-projection conjecture}
\author{Autonomous proof packet for arXiv:1909.00391}
\date{11 August 2026}

\begin{document}
\maketitle

\section*{Problem and result}

Conjecture 6.1 of Arhancet asks, for $1<p<\infty$, $p\ne2$, whether the
positively contractively complemented subspaces of Haagerup noncommutative
$L^p$-spaces are exactly the nonassociative $L^p$-spaces of $JW^*$-algebras
with faithful normal states.  The exact source statement is reproduced below.

\begin{center}
\includegraphics[width=.91\textwidth]{figures/open_conjecture_crop.png}
\end{center}

The later companion paper \cite{Arhancet2023} defines
$L^p(\F,\tau)=(\F,\F_*)_{1/p}$ using the trace-induced compatibility and
proves the embedding implication for a single tracial $JW^*$-factor with
separable predual.  We extend that implication across a substantial class of
nonfactorial algebras.

\begin{theorem}[piecewise-homogeneous tracial embedding]
Let $1\le p<\infty$.  For each $j$ in a finite or countable set $I$, let
$(\Omega_j,\mu_j)$ be a finite measure space and let $\F_j$ be a $JW^*$-factor
with separable predual and normalized normal faithful trace $\tau_j$.  Put
\[
 \J=\bigoplus_{j\in I}^{\ell^\infty}
 L^\infty(\Omega_j,\mu_j;\F_j)
\]
and equip it with the normal finite faithful trace
\[
 \tau(x)=\sum_{j\in I}a_j\int_{\Omega_j}\tau_j(x_j(\omega))\,d\mu_j(\omega),
 \qquad a_j>0,
\]
where the weights make $\tau(1)<\infty$.  Then the nonassociative space
$L^p(\J,\tau)$ is isometric to the range of a positive contractive projection
on $L^p(\M,\sigma)$ for a finite von Neumann algebra $(\M,\sigma)$.
\end{theorem}

Thus the embedding direction of Conjecture 6.1 holds for all such tracial
$\J$.  This includes countable atomic centers ($\Omega_j$ singletons) and
homogeneous diffuse centers.

\section*{Proof}

Fix $j$.  By Theorem 4.7 of \cite{Arhancet2023}, there are a finite von
Neumann algebra $(\N_j,s_j)$ and a positive contractive projection
\[
 R_j:L^p(\N_j,s_j)\longrightarrow L^p(\N_j,s_j)
\]
whose range is isometric to $L^p(\F_j,\tau_j)$.  In the construction of that
theorem the projection is induced by a trace-preserving normal positive
projection at the algebra level; equivalently, it may be applied pointwise.

Set
\[
 \M_j=L^\infty(\Omega_j,\mu_j)\,\bar\otimes\,\N_j,
 \qquad
 \sigma_j=\mu_j\otimes s_j.
\]
The induced map
\[
 \widetilde R_j f(\omega)=R_j(f(\omega))
\]
is a positive contractive projection on
$L^p(\M_j,\sigma_j)=L^p(\Omega_j;L^p(\N_j,s_j))$.  Its range is
\[
 L^p(\Omega_j;L^p(\F_j,\tau_j)).                                      \tag{1}
\]
On the other hand, exact vector-valued complex interpolation gives
\begin{align*}
 L^p(L^\infty(\Omega_j;\F_j),\mu_j\otimes\tau_j)
 &=\big(L^\infty(\Omega_j;\F_j),L^1(\Omega_j;(\F_j)_*)\big)_{1/p}\\
 &=L^p\big(\Omega_j;(\F_j,(\F_j)_*)_{1/p}\big)\\
 &=L^p(\Omega_j;L^p(\F_j,\tau_j))                                    \tag{2}
\end{align*}
isometrically.  Hence (1) is exactly the desired homogeneous
nonassociative $L^p$-space, not merely an isomorphic copy with uncontrolled
constants.

Now take the finite-tracial von Neumann algebraic direct sum
$\M=\bigoplus_{j\in I}^{\ell^\infty}\M_j$ and give its $j$th summand the
central trace weight $a_j$.  Standard central decomposition yields
\[
 L^p(\M,\sigma)=
 \left(\bigoplus_{j\in I}L^p(\M_j,\sigma_j)\right)_{\ell^p(a_j)}.
\]
The coordinatewise map $R=\bigoplus_j\widetilde R_j$ is positive and
idempotent, and
\[
 \|Rx\|_p^p=\sum_j a_j\|\widetilde R_jx_j\|_p^p
 \le \sum_j a_j\|x_j\|_p^p=\|x\|_p^p.
\]
Its range is the same weighted $\ell^p$-sum of the spaces in (2).  Applying
the identical central-sum interpolation formula to $(\J,\J_*)$ identifies
that range isometrically with $L^p(\J,\tau)$.  This proves the theorem. \qed

\section*{Proof intuition}

There are two exact stability principles.  A factor solution survives adding
a scalar measure variable because complex interpolation commutes with
$L^p(\Omega;\cdot)$.  It also survives a countable central sum because both
the ambient and nonassociative $L^p$ norms become the same weighted
$\ell^p$ norm.  Positivity and contractivity can therefore be checked one
fiber at a time.

The point is not just closure under finite direct sums: diffuse homogeneous
centers are included as well.  What remains outside the theorem is a genuinely
varying measurable field of non-isomorphic factors, where one needs a
measurable choice of ambient von Neumann algebra and factor projection.

\section*{Scope and remaining obstruction}

This does not prove Conjecture 6.1.  The companion paper itself leaves the
arbitrary tracial algebra as Conjecture 5.3 and notes that the direct-integral
formalism is unclear.  Our argument removes this issue only for countably
piecewise-constant fields.  For nontracial states there is a further modular
obstruction: a state-preserving Jordan conditional expectation need not be
selfadjoint, and state-independence of $(\J,\J_*)_{1/p}$ is itself open.
Finally, the general forward implication still requires identifying the range
of the lifted positive projection with the appropriate interpolation space.

\begin{thebibliography}{9}
\bibitem{Arhancet2019}
C. Arhancet, \emph{Positive contractive projections on noncommutative
$L^p$-spaces and nonassociative $L^p$-spaces}, arXiv:1909.00391v8.

\bibitem{Arhancet2023}
C. Arhancet, \emph{Nonassociative $L^p$-spaces and embeddings in
noncommutative $L^p$-spaces}, arXiv:2307.04452v3, especially Theorem 4.7 and
Conjecture 5.3.
\end{thebibliography}

\end{document}
